\section{Introduction}
\label{sec:introduction}

As large language models (LLMs) scale to billions of parameters, full fine-tuning becomes increasingly expensive. The dominant alternative is parameter-efficient adaptation, in which only a small set of additional parameters is introduced \citep{hu2022lora, houlsby2019adapter}. Among these, \textit{soft prompts} insert learnable continuous embeddings outside the vocabulary into the input, steering model behavior through the attention mechanism \citep{li2021prefix, lester2021power}. The same paradigm has been broadly adopted to enhance mathematical reasoning in recent work \citep{kang2025model, xu2025softcot, ye2025thinking, hao2024training, zhang2025memgen}, where the common structure is to \emph{optimize} continuous embeddings outside the vocabulary to improve downstream accuracy.

There is, however, a clue that the effect does not come from optimization alone: injecting untrained random continuous embeddings already shifts the output distribution. For example, applying a chat template to a base model that was not fine-tuned for chat causes a format mismatch that degrades reasoning accuracy, but appending random embeddings recovers a substantial portion of that loss. If simple noise were merely perturbing the model, accuracy should worsen instead. This suggests that the effect arises not from \emph{learned content} but from some property of the \emph{injection itself}---a question prior evaluations, focused on end-task accuracy, have not separated from learned content.

We investigate Random Soft Prompts (RSPs), continuous vectors injected without any training. Each RSP entry is sampled from an isotropic Gaussian fitted to the entrywise mean and variance of the pretrained embedding table. The role of randomness is not to imitate the embedding distribution; it is to ensure that repeated rollouts receive \emph{independent} perturbations. We answer two questions. First, why can one sampled RSP perturb early next-token decisions and then fade? An exact attention decomposition and two envelope bounds show that the same attention mass both injects and attenuates the perturbation, and that cache growth shrinks the admissible mass late in decoding. Second, why do fresh RSP draws accumulate across trials? Under a fixed variance budget, isotropic perturbations are direction-agnostic, and the Gaussian law additionally assigns positive probability to every nonempty open set. Independent resampling can therefore reopen local top-$K$ entry regions across rollouts, which becomes Pass@$N$ growth once a separate task-side success assumption holds. Finally, we discuss what this diversity implies for GRPO \citep{shao2024deepseekmath} training.

Our contributions are as follows.
\begin{itemize}
  \item We empirically analyze how RSPs affect LLM generation: RSPs raise early-token entropy and stabilize back to baseline at later steps, and combined with temperature sampling produce more diverse reasoning trajectories.
  \item We explain why untrained RSPs can matter at the transformer level. A single attention-side scalar controls both the size of the injected perturbation and its late-stage decay, and the centered isotropic prompt law neither induces a fixed first-order drift nor excludes any local open branch region.
  \item We discuss what RSP-induced diversity implies for GRPO training. RSP is a \emph{rank-changing} perturbation while temperature is \emph{rank-preserving}---this difference makes the two methods provide orthogonal exploration gains.
\end{itemize}
